\chapter{Wie Menschen wirklich entscheiden}

\section{Dieselbe Lage, zwei Formulierungen}

\begin{merksatz}
\boxtitel{Der Rahmen entscheidet mit}
Eine Seuche bedroht $600$ Menschen. Zwei Programme stehen zur Wahl.

\begin{center}
{\small
\renewcommand{\arraystretch}{1.5}
\begin{tabular}{@{}P{1.4cm}P{6.8cm}@{}}
\toprule
\multicolumn{2}{@{}l}{\color{akzent}\bfseries Fassung 1 --- Überlebende}\\
$A$ & $200$ werden gerettet.\\
$B$ & Mit $\tfrac13$ werden alle gerettet, mit $\tfrac23$ niemand.\\
\midrule
\multicolumn{2}{@{}l}{\color{akzent}\bfseries Fassung 2 --- Tote}\\
$C$ & $400$ sterben.\\
$D$ & Mit $\tfrac13$ stirbt niemand, mit $\tfrac23$ sterben alle.\\
\bottomrule
\end{tabular}}
\end{center}

\begin{center}
\begin{minipage}{0.9\textwidth}
\textbf{$A$ und $C$ beschreiben dasselbe. $B$ und $D$ auch. In Befragungen wählt die Mehrheit $A$ --- und $D$.}
\end{minipage}
\end{center}

\begin{center}
\begin{minipage}{0.88\textwidth}
\kursiv{\textbf{Wer $A$ vorzieht, ist im Gewinnbereich risikoscheu; wer $D$ vorzieht, im Verlustbereich risikofreudig.} Der Unterschied liegt nicht in der Lage, sondern in dem Punkt, von dem aus gezählt wird. Nach Satz 2.1 dürfte das keine Rolle spielen --- die Nutzenfunktion bewertet Endzustände, nicht Veränderungen.}
\end{minipage}
\end{center}
\end{merksatz}

\section{Was man stattdessen beobachtet}

\begin{merksatz}
\boxtitel{Drei Abweichungen}
\begin{center}
{\small
\renewcommand{\arraystretch}{1.55}
\begin{tabular}{@{}P{2.8cm}P{5.8cm}@{}}
\toprule
\textbf{Bezugspunkt} & Bewertet werden Gewinne und Verluste gegenüber einem Ausgangspunkt, nicht Endzustände.\\
\textbf{Verlustscheu} & Ein Verlust wiegt etwa doppelt so schwer wie ein gleich großer Gewinn.\\
\textbf{Gewichtete Wahrscheinlichkeit} & Kleine Wahrscheinlichkeiten werden überschätzt, mittlere und große unterschätzt.\\
\bottomrule
\end{tabular}}
\end{center}

\schritt{Die dritte Zeile in Zahlen}
\begin{center}
{\small
\renewcommand{\arraystretch}{1.48}
\begin{tabular}{@{}P{2.4cm}P{2.6cm}P{2.4cm}@{}}
\toprule
\textbf{wahr} & \textbf{gewichtet} & \\
\midrule
$0{,}01$ & $\approx 0{,}05$ & überschätzt\\
$0{,}10$ & $\approx 0{,}19$ & \\
$0{,}50$ & $\approx 0{,}42$ & unterschätzt\\
$0{,}99$ & $\approx 0{,}95$ & \\
\bottomrule
\end{tabular}}
\end{center}

\begin{center}
\begin{minipage}{0.88\textwidth}
\kursiv{\textbf{Das erklärt Lotterie und Versicherung zugleich}: Beide beruhen auf der Überschätzung kleiner Wahrscheinlichkeiten, einmal für den Gewinn, einmal für den Schaden. In der Medizin erklärt es, warum seltene, benannte Nebenwirkungen eine Behandlung kippen können, obwohl ihr Erwartungsbeitrag winzig ist --- und warum Patienten häufige, unauffällige Risiken achselzuckend hinnehmen.}
\end{minipage}
\end{center}
\end{merksatz}

\section{Was das für die Aufklärung bedeutet}

\begin{merksatz}
\boxtitel{Eine Zahl, vier Darstellungen}
Die Gerinnungshemmung senkt das Schlaganfallrisiko von sechs auf zwei Prozent im Jahr.

\begin{center}
{\small
\renewcommand{\arraystretch}{1.5}
\begin{tabular}{@{}P{3.2cm}P{2.4cm}P{2.6cm}@{}}
\toprule
\textbf{Darstellung} & \textbf{Zahl} & \textbf{Wirkung}\\
\midrule
relative Senkung & $67\,\%$ & wirkt groß\\
absolute Senkung & $4$ Punkte & wirkt mäßig\\
Zahl der Behandelten & $25$ & wirkt klein\\
Überlebende ohne Ereignis & $94 \to 98$ & wirkt klein\\
\bottomrule
\end{tabular}}
\end{center}

\begin{center}
\begin{minipage}{0.9\textwidth}
\textbf{Alle vier Angaben stimmen. Die Zustimmungsraten unterscheiden sich in Befragungen um das Zwei- bis Dreifache.}
\end{minipage}
\end{center}

\begin{center}
\begin{minipage}{0.88\textwidth}
\kursiv{\textbf{Es gibt keine rahmenfreie Darstellung} --- auch das Weglassen ist eine Wahl. Man kann nur mehrere Darstellungen nebeneinanderstellen und offenlegen, dass sie dasselbe bedeuten. Ärzte sind gegen diese Wirkungen nicht immun: In Befragungen empfahlen sie dieselbe Behandlung häufiger, wenn sie als Überlebensrate statt als Sterberate beschrieben war.}
\end{minipage}
\end{center}
\end{merksatz}

\section{Der Bruchfall}

\begin{warnung}
\boxtitelr{Zwei Fehlschlüsse}
\begin{center}
{\small
\renewcommand{\arraystretch}{1.52}
\begin{tabular}{@{}P{2.8cm}P{5.8cm}@{}}
\toprule
\textbf{Abweichung ist Irrtum} & Nicht jede Abweichung ist ein Fehler. Wer den Verlust fürchtet, mag gute Gründe haben --- Kapitel~2 verlangt nur Widerspruchsfreiheit, nicht Gleichmut.\\
\textbf{Wer weiß, was gut ist} & Aus der Beobachtung, dass Menschen anders entscheiden, folgt nicht, dass jemand anderes für sie entscheiden soll.\\
\bottomrule
\end{tabular}}
\end{center}

\begin{center}
\begin{minipage}{0.88\textwidth}
\kursiv{\textbf{Die zweite Zeile trifft die Debatte um das \enquote{Anstupsen}}: Wenn jede Darstellung rahmt, dann rahmt auch die bewusst gewählte --- und die Frage ist nicht mehr, ob gerahmt wird, sondern wer es tut und wozu. In der Medizin ist die Antwort einfacher als in der Politik: Der Rahmen soll die Bewertung des Patienten hervorbringen, nicht die des Arztes.}
\end{minipage}
\end{center}
\end{warnung}

\section{Übungsaufgaben}

\begin{aufgabe}
\boxtitel{Aufgabe 12.1 --- Verlustscheu rechnen}
Jemand bewertet Gewinne mit $v(x) = x$ und Verluste mit $v(x) = 2{,}25\,x$. Nimmt er eine Wette \enquote{mit je einhalb $+100$ oder $-80$}?

\loesung
\schritt{Schritt 1 --- Der Erwartungswert}
\[ 0{,}5\cdot100 - 0{,}5\cdot80 = +10 \]

\schritt{Schritt 2 --- Der bewertete Wert}
\[ 0{,}5\cdot100 - 0{,}5\cdot2{,}25\cdot80 = 50 - 90 = -40 \]

\begin{center}
\begin{minipage}{0.9\textwidth}
\textbf{Er lehnt ab --- obwohl die Wette im Mittel zehn einbringt.}
\end{minipage}
\end{center}

\schritt{Schritt 3 --- Ab wann er annimmt}
\[ 0{,}5\,G = 0{,}5\cdot2{,}25\cdot80 \quad\Longrightarrow\quad G = 180 \]

\begin{center}
{\small
\renewcommand{\arraystretch}{1.48}
\begin{tabular}{@{}P{2.4cm}P{2.4cm}P{2.4cm}@{}}
\toprule
\textbf{Verlust} & \textbf{nötiger Gewinn} & \textbf{Verhältnis}\\
\midrule
$80$ & $180$ & $2{,}25$\\
$1000$ & $2250$ & $2{,}25$\\
\bottomrule
\end{tabular}}
\end{center}

\schritt{Schritt 4 --- In der Klinik}
\begin{center}
\begin{minipage}{0.88\textwidth}
\kursiv{\textbf{Eine Behandlung, die einen Nutzen bringt und ein gleich großes Risiko trägt, wird abgelehnt} --- und zwar umso deutlicher, je klarer das Risiko benannt ist. Die Ablehnung einer Impfung mit kleinem, benanntem Nebenwirkungsrisiko gegenüber einer größeren, aber diffusen Krankheitsgefahr folgt genau diesem Muster: benannter Verlust gegen unbestimmten Gewinn.}
\end{minipage}
\end{center}
\end{aufgabe}

\begin{aufgabe}
\boxtitel{Aufgabe 12.2 --- Den Bezugspunkt verschieben}
Ein Patient erfährt zunächst, seine Lebenserwartung betrage fünf Jahre; später, sie betrage acht. Eine Behandlung verspricht zehn. Wie wirkt die Reihenfolge?

\loesung
\schritt{Schritt 1 --- Vom ersten Bezugspunkt aus}
\[ +5 \ \text{Jahre Gewinn} \]

\schritt{Schritt 2 --- Vom zweiten aus}
\[ +2 \ \text{Jahre Gewinn} \]

\schritt{Schritt 3 --- Die Bewertung}
\begin{center}
{\small
\renewcommand{\arraystretch}{1.48}
\begin{tabular}{@{}P{2.6cm}P{2.2cm}P{2.6cm}@{}}
\toprule
\textbf{Bezugspunkt} & \textbf{Gewinn} & \textbf{Endzustand}\\
\midrule
fünf Jahre & $+5$ & $10$\\
acht Jahre & $+2$ & $10$\\
\bottomrule
\end{tabular}}
\end{center}

\begin{center}
\begin{minipage}{0.9\textwidth}
\textbf{Der Endzustand ist derselbe. Nach Satz 2.1 muss die Bewertung gleich sein; tatsächlich ist sie es nicht.}
\end{minipage}
\end{center}

\schritt{Schritt 4 --- Die praktische Folge}
\begin{center}
\begin{minipage}{0.88\textwidth}
\kursiv{\textbf{Wer zuerst eine düstere Prognose nennt und dann eine Behandlung anbietet, erzeugt ein Gefühl des Gewinns} --- wer zuerst Hoffnung macht und dann einschränkt, eines des Verlusts. Beides ist dieselbe Mitteilung. Die Reihenfolge ärztlicher Aufklärung ist deshalb keine Nebensache, und sie lässt sich nicht neutral gestalten; man kann sie nur bewusst wählen und begründen.}
\end{minipage}
\end{center}
\end{aufgabe}

\begin{aufgabe}
\boxtitel{Aufgabe 12.3 --- Die vier Darstellungen prüfen}
Rechnen Sie die vier Angaben aus Abschnitt 12.3 nach.

\loesung
\schritt{Schritt 1 --- Relative Senkung}
\[ \frac{0{,}06-0{,}02}{0{,}06} = \frac{0{,}04}{0{,}06} = 0{,}667 = 66{,}7\,\% \]

\schritt{Schritt 2 --- Absolute Senkung}
\[ 0{,}06 - 0{,}02 = 0{,}04 = 4 \ \text{Punkte} \]

\schritt{Schritt 3 --- Zahl der notwendigen Behandlungen}
\[ \frac{1}{0{,}04} = 25 \]

\schritt{Schritt 4 --- Ereignisfreie je hundert}
\[ 94 \to 98 \]

\schritt{Schritt 5 --- Die Probe}
\begin{center}
{\small
\renewcommand{\arraystretch}{1.48}
\begin{tabular}{@{}P{3.4cm}P{2.4cm}@{}}
\toprule
$25$ Behandelte, ohne & $1{,}5$ Ereignisse\\
$25$ Behandelte, mit & $0{,}5$ Ereignisse\\
\midrule
verhindert & $\mathbf{1{,}0}$ \checkmark\\
\bottomrule
\end{tabular}}
\end{center}

\begin{center}
\begin{minipage}{0.88\textwidth}
\kursiv{\textbf{Die Zahl der notwendigen Behandlungen ist der Kehrwert der absoluten Senkung} --- und die einzige der vier Angaben, die die Größenordnung unmittelbar anschaulich macht: Vierundzwanzig von fünfundzwanzig Behandelten haben von der Behandlung nichts, tragen aber ihr Blutungsrisiko. Genau deshalb wird sie in Werbematerialien selten genannt.}
\end{minipage}
\end{center}
\end{aufgabe}

\begin{geschichte}
\boxtitelg{Zwei Psychologen und ein Nobelpreis}
\Person{Daniel Kahneman} und \Person{Amos Tversky} veröffentlichten 1979 die Neue Erwartungstheorie --- eine Beschreibung dessen, wie Menschen tatsächlich wählen.

\begin{center}
\begin{minipage}{0.9\textwidth}
\kursiv{Der Aufsatz erschien in \kursiv{Econometrica}, der strengsten Zeitschrift der Wirtschaftswissenschaft, obwohl beide Verfasser Psychologen waren. \textbf{Er ist bis heute einer der meistzitierten Aufsätze der Sozialwissenschaften überhaupt.} Kahneman erhielt 2002 den Wirtschaftsnobelpreis; Tversky war 1996 gestorben, und der Preis wird nicht postum verliehen.}
\end{minipage}
\end{center}

\schritt{Das Seuchenbeispiel}
\begin{center}
\begin{minipage}{0.9\textwidth}
\kursiv{Es stammt aus einer Arbeit von 1981 und wurde in Dutzenden Fächern nachgestellt, auch mit Ärzten und mit Studienteilnehmern, die beide Fassungen nacheinander sahen. \textbf{Selbst wer den Widerspruch bemerkt, ändert seine Antwort häufig nicht} --- die Wirkung verschwindet nicht durch Einsicht.}
\end{minipage}
\end{center}

\schritt{Und die Gegenrede}
\begin{center}
\begin{minipage}{0.9\textwidth}
\kursiv{\Person{Gerd Gigerenzer} wandte seit den 1990er Jahren ein, viele dieser Abweichungen verschwänden, wenn man die Aufgaben in natürlichen Häufigkeiten statt in Wahrscheinlichkeiten stelle --- \enquote{zehn von tausend} statt \enquote{ein Prozent}. \textbf{Damit wurde aus einem Befund über menschliche Grenzen eine Frage der Darstellung}, und die medizinische Aufklärung hat daraus mehr gelernt als aus der Theorie selbst.}
\end{minipage}
\end{center}
\end{geschichte}

\begin{forschung}
\boxtitelp{Beschreibende Entscheidungsforschung heute}
\begin{center}
{\small
\renewcommand{\arraystretch}{1.45}
\begin{tabular}{@{}P{3.4cm}P{4.8cm}@{}}
\toprule
\textbf{Natürliche Häufig\-keiten} & Darstellungen, die Rechenfehler vermeiden helfen.\\
\textbf{Risikokommuni\-kation} & Wie man Nutzen und Schaden so zeigt, dass beides gleich sichtbar ist.\\
\textbf{Wiederholbarkeit} & Viele klassische Befunde sind kleiner als gedacht.\\
\bottomrule
\end{tabular}}
\end{center}

\schritt{Zur dritten Zeile}
\begin{center}
\begin{minipage}{0.88\textwidth}
\kursiv{Die Wiederholungskrise der Psychologie hat auch dieses Gebiet erfasst. \textbf{Die großen Effekte --- Rahmung, Verlustscheu, Überschätzung kleiner Wahrscheinlichkeiten --- haben sich gehalten}, viele feinere nicht. Für die Medizin genügt der harte Kern: Die Darstellung wirkt, und zwar stark genug, um Entscheidungen zu kippen.}
\end{minipage}
\end{center}
\end{forschung}

\begin{merksatz}
\boxtitel{Was Sie jetzt können}
\begin{center}
\renewcommand{\arraystretch}{1.5}
\begin{tabular}{@{}P{3.0cm}P{5.6cm}@{}}
\toprule
\textbf{Rahmung} & dieselbe Lage, andere Wahl\\
\textbf{Bezugspunkt} & Gewinne und Verluste statt Endzustände\\
\textbf{Verlustscheu} & Faktor etwa $2{,}25$\\
\textbf{Gewichtung} & kleine Wahrscheinlichkeiten überschätzt\\
\textbf{Vier Darstellungen} & alle richtig, alle verschieden wirksam\\
\bottomrule
\end{tabular}
\end{center}

\begin{center}
\begin{minipage}{0.88\textwidth}
\kursiv{Damit endet Teil~IV. Bisher entschied einer für sich. Teil~V fragt, was geschieht, wenn mehrere gemeinsam entscheiden --- und wenn für viele zugleich entschieden wird.}
\end{minipage}
\end{center}
\end{merksatz}
